% !PRINTED 309
\weberpage{309}
Hierbei ist noch folgende Bemerkung von Interesse. Die Resolventen $p$ten Grades, deren Wurzeln die Größen (9), (10), (11) sind, enthalten nach §\,78 in ihren Koeffizienten noch $\sqrt{\pm p}$.

Die Gruppe $\mathfrak{L}_0$ wird aber zur Gruppe $\mathfrak{L}$ erweitert, wenn wir \gesperrt{eine} lineare Substitution von der Form §\,78, (13) hinzufügen, in der $ad - bc$ quadratischer Nichtrest ist, also etwa: % Source scan clearly prints ad - bc.
\[
\begin{aligned}
\text{für } p&=5 & &(z,\,2z)\\
\text{für } p&=7,\,11 & &(z,\,-z).
\end{aligned}
\]
Durch Anwendung dieser Substitution geht aber der Wert (9)
\[
V = (v_\infty - v_0)(v_1 - v_4)(v_2 - v_3)
\]
in den entgegengesetzten über, und daraus folgt, daß, wenn $p=5$ ist, $V^2$ einer Gleichung 5ten Grades mit \gesperrt{rationalen} Zahlenkoeffizienten genügt. Daraus schließt man, daß in der Gleichung für $V$ die Koeffizienten der ungeraden Potenzen den Faktor $\sqrt5$ haben, während die anderen rational sind, oder daß man für die Unbekannte $\sqrt5\,V$ eine Gleichung mit \gesperrt{rationalen Koeffizienten} erhält.

Für $p=7,\,11$ gehen die den beiden Vorzeichen in (10) oder (11) entsprechenden Ausdrücke durch die Vertauschung $(z,-z)$ ineinander über. Die Resolventen 7ten oder 11ten Grades, denen einer der beiden Ausdrücke (10) bzw. (11) genügt, gehen also durch Änderung des Vorzeichens von $\sqrt{-7}$, $\sqrt{-11}$ ineinander über.

§\,83. \gesperrt{Verschiedene Resolventen 5ten Grades für den 5ten Transformationsgrad.}

Bei der Bildung der Resolventen $p$ten Grades beschränken wir uns hier auf den Fall $p=5$.

Diesen Resolventen kann man sehr mannigfaltige Formen geben, indem man nicht nur in der Funktion $V$, (9) des vorigen Paragraphen für die Größen $v$ die Wurzeln einer \gesperrt{beliebigen} Transformationsgleichung wählen, sondern auch $V$ durch mancherlei andere Funktionen ersetzen kann, etwa durch
\[
(v_\infty + v_0)(v_1+v_4)(v_2+v_3),
\]
oder durch
\[
v_\infty v_0 + v_1 v_4 + v_2 v_3.
\]
